\section{Character-sensitive transfer with a fixed multiplier ideal}
\label{inj:character-transfer}

The next conversion allows a character shift of the anticanonical
linearization.  Its proof uses a common field of invariant section
ratios and a singular metric with positive curvature in the base
directions.  The boundary correction on the base has exceptional
pullback, and its multiplier ideal stays fixed while the exponent
varies.

\begin{proof}[Proof of Theorem~\ref{inj:transfer}]
\smallskip\noindent\textit{The common ratio field.} Put $n=\dim X$ and write $N_i=M+iP$.  For each $i$, ratios of invariant sections of
the same positive multiple of $N_i$ form a field $Q_i$.
For $p<q<b$ in $I$, the identity
\[
 (b-p)N_q=(b-q)N_p+(q-p)N_b
\]
holds with linearizations.  Multiplication by suitable powers of
the fixed $s_p,s_b$ puts any ratio for $N_p$ or $N_b$ into $Q_q$.
This retains the ratio itself: for a ratio in degree $aN_p$, multiply
its numerator and denominator by
$s_p^{a(b-q-1)}s_b^{a(q-p)}$, and use the analogous formula for
a ratio in $aN_b$.  Taking a further index beyond any prescribed
two interior indices proves that all $Q_i$ with $i>\min I$ coincide;
denote this field by $Q$.

Fix the first two indices $p<q$.  For $b>q$ the ratio
\begin{equation}\label{inj:ratio-field-identity}
 \frac{s_q^{b-p}}{s_p^{b-q}s_b^{q-p}}
\end{equation}
belongs to $Q$.  If $\operatorname{trdeg}_{\mathbb C}Q=0$,
it is constant.  The section divisors are affine in $i$, and their
slope is effective because the indices are unbounded.  Thus
$s_q/s_p$ is a regular invariant section of $(q-p)P$.

Otherwise put $e=\operatorname{trdeg}Q>0$.  Choose $e$ independent
ratios to define a map to $\mathbb P^e$, resolve its graph and take
Stein factorization. The Stein field is the relative algebraic closure
in $\mathbb C(X)$ of the field of the chosen ratios, a finite extension
of that field. It contains $Q$, since $Q$ is algebraic over the chosen
ratio field, and is relatively algebraically closed in $\mathbb C(X)$.
The connected torus acts
trivially on the Stein base because its orbits lie in finite fibers
over a base with trivial action.  As before take a smooth projective
base $Y$, a normal equidimensional $W\to Y$ with a birational
morphism $W\to X$, and an equivariant smooth resolution $V\to W$,
with maps $\mu:W\to X$, $f_W:W\to Y$, and total maps $\pi:V\to X$ and $f:V\to Y$. Put $P_W=\mu^*P$.  The map $f$ has
connected fibers by the Stein construction and normality of $Y$.
Write $L'=\pi^*L$ and
$P'=\pi^*P$.

For each interior index $i$ and every divisor $H'$ on $Y$,
\begin{equation}\label{inj:ratio-domination}
 H^0(V,a\pi^*N_i-f^*H')^T\ne0
 \quad\text{for some integer }a>0.
\end{equation}
Indeed the chosen independent ratios belong to $Q_i$.  Put their
numerators and denominator in one common degree by products of
invariant sections.  On the resolved graph they have a common
invariant factor whose line is that multiple of $N_i$ minus the
pullback of $\mathcal O_{\mathbb P^e}(1)$.  The latter line is
big on $Y$, so a sufficiently large multiple dominates $H'$.

\smallskip\noindent\textit{A regular relative section.} We next remove the vertical poles while preserving the character. Let $D_i$ be the section divisor on $W$.  Since
\eqref{inj:ratio-field-identity} comes from $Y$, its horizontal
divisor is zero.  The affine relations show that the horizontal
part of $(D_q-D_p)/(q-p)$ is effective.  Normalize the vertical
part by subtracting the minimum coefficient divided by pullback
multiplicity above each base prime, and let $B$ be the resulting
rational base divisor.  There is an integer $k>0$, divisible by
$q-p$ and clearing the coefficients of $B$, for which
\begin{equation}\label{inj:gamma}
 \gamma=(s_q/s_p)^{k/(q-p)}(f_W^*1_{kB})^{-1}
 \in H^0(W,kP_W-f_W^*(kB))^T
\end{equation}
is regular.  Normality checks regularity in codimension one, and
equidimensionality accounts for every vertical prime.  Above every
base prime $\gamma$ is nonzero at the general point of at least
one component.  We also write $\gamma$ for its pullback to $V$.

\smallskip\noindent\textit{All adjoint ranks, including zero.} The ratio field must also control the fiber integrals. Let $E=K_{V/X}\ge0$.  Its Jacobian section $1_E$ is invariant
for the natural linearization on $L'$, irrespective of the character
shift in $P'$.  For all $j\ge0$,
\begin{equation}\label{inj:transfer-rank}
 \operatorname{rank}
 \bigl(f_*\mathcal O_V(E+jkP')\bigr)^T=1.
\end{equation}
The section $1_E\gamma^j$ gives the lower bound over the generic
fiber.  If the rank were larger, globalize two invariant sections
$u_1,u_2$ independent over $\mathbb C(Y)$ after a twist $f^*(aH)$,
where $H$ is ample.  Choose $i,b\in I$ with $b>i+jk$.  Multiply
$u_1^k$ and $u_1^{k-1}u_2$ by
$s_i^k\gamma^{b-i-jk}$.  Both results are sections of
\[
 kE+k\pi^*N_b+
 f^*\bigl(kaH-(b-i-jk)kB\bigr),
\]
with ratio $u_2/u_1$.  Apply \eqref{inj:ratio-domination} at
index $b$ to cancel this base twist by the same invariant multiplier
in the numerator and denominator.  Removing the exceptional twist
$kE$ and pushing to $X$ leaves two invariant sections of an equal
multiple of $N_b$.  Normality extends across the image of the
exceptional locus.  Their ratio must be in $Q\subset\mathbb C(Y)$,
contrary to independence.  This proves \eqref{inj:transfer-rank}.

\smallskip\noindent\textit{The bounded metric and the horizontally positive metric.} We now construct the two metrics needed for interpolation. Average a smooth semipositive metric $h_0$ on $L$ over the maximal
compact torus.  It is also invariant for $P$, because the two
linearizations differ by a character, unitary on that compact torus.
Equation \eqref{inj:ratio-domination} and pseudoeffectivity of $-M$
give $L'-\eta f^*H$ pseudoeffective for some $\eta>0$ and ample
$H$.  Taking a positive current, averaging, and mixing a small
positive fraction of its potential with $h_0$ gives an invariant
singular metric $h$ on $L'$ with
\begin{equation}\label{inj:strict-integrable}
 i\Theta_h\ge\delta f^*\omega_Y,
 \qquad \mathcal I(h)=\mathcal O_V,
 \qquad\delta>0.
\end{equation}
Skoda integrability and compactness give the last equality.  The
metric is bounded below by a positive constant times $h_0$ on
compact coordinate patches: its potential relative to the smooth
metric is locally bounded above.

Apply Proposition~\ref{har:prop:maximum} with the normal model $W$, the chosen linearization $P$, and the relative section $\gamma$ of $kP_W-f_W^*(kB)$. Equivalently use the rational section $(s_q/s_p)^{k/(q-p)}$ before normalization. Equation~\eqref{inj:transfer-rank} gives every invariant adjoint rank, the Jacobian has its natural invariant action, and the averaged smooth metric is invariant also for $P$. The vertical minima defining $kB$ are exactly those in that proposition. Its smooth coefficient is $L'+jkP'$, ordinary $(1+jk)L'$, so the orthogonal quotient argument applies without any singular-subbundle claim. We obtain a continuous, in particular locally bounded, semipositive metric on $kB$.

\smallskip\noindent\textit{A fixed ideal on the base.} The continuous base metric just obtained controls the varying power. The singular degree-zero metric will supply the positive constant term.
For $j=0$ use the singular coefficient metric $h$.
The singular adjoint direct-image theorem
\cite[Setup 3.2.1 and Theorem 3.3.4]{PT2018} applies over the
smooth, locally free base-change locus.  Its generic multiplier-ideal
inclusion is an isomorphism by \eqref{inj:strict-integrable}; fiber
integrability is an almost-everywhere assertion by Fubini.
The invariant summand is an orthogonal holomorphic quotient for
the compact-invariant metric.  Thus its line metric has
semipositive curvature.  Subtract a local potential for
$\delta\omega_Y$ from the coefficient weight to see that its
curvature is in fact at least $\delta\omega_Y$.

In a frame $\xi$ of $-K_Y$, denote the squared norm of
$1_Ef^*\xi$ by $F_\xi(y)$.  With base volume in the dual frame,
this is the pushforward density of the integrable absolute adjoint
form $1_E$ with metric $h$.  Hence $F_\xi$ is locally integrable
on the whole base chart.  We next extend its negative logarithm
with an explicit exceptional correction.

\smallskip\noindent\textit{Boundary extension with exceptional support.} For a prime boundary divisor $P$ choose a prime $Q\subset V$
mapping onto it, of pullback multiplicity $l$ and Jacobian order
$a=\operatorname{ord}_Q(E)$.  At general points write
$f=(z_1^l,z_2,\ldots,z_e)$.  Integrating on one coordinate patch,
using $h\ge c h_0$, gives almost everywhere
\begin{equation}\label{inj:density-bound}
 F_\xi(y)\ge c'|y_1|^{2((a+1)/l-1)}.
\end{equation}
The exponent includes the transverse Jacobian factor
$|z_1|^{-2(l-1)}$. If a nonexceptional component dominates $P$, choose
it for the estimate. Then $a=0$ and the coefficient of $P$ in the
correction is zero. Otherwise every component dominating $P$ is
exceptional; choose a nonnegative integer coefficient at least
$(a+1)/l-1$. Summing
over the finitely many boundary primes gives $D_*\ge0$ with
$f^*D_*$ exceptional over $X$.

For the last assertion one must also check primes mapping to
codimension at least two on $Y$.  The equidimensional normal model
$W$ has relative dimension $n-e$, so for $Z\subset Y$ of
codimension at least two,
\[
 \dim f_W^{-1}(Z)\le\dim Z+n-e\le n-2.
\]
Any prime of $V$ over $Z$ therefore maps to codimension at least
two on $W$ and on $X$. Every component dominating a base prime in
$\Supp D_*$ is exceptional by the rule defining that support. These two checks
account for every component of $f^*D_*$.

On $-K_Y+D_*$ define the metric by
$|\xi1_{D_*}|^2=F_\xi$ on the dense open.
Its ordinary-frame weight is
$-\log F_\xi+\log|g_{D_*}|^2$.  Estimate
\eqref{inj:density-bound} bounds this weight above at general
boundary points. Subtract a local smooth potential for
$\delta\omega_Y$. The resulting plurisubharmonic weight extends
across each boundary divisor away from a codimension-two analytic
subset. Hartogs extension for plurisubharmonic functions then fills
that subset. Restoring the smooth potential gives a metric on all of
$Y$ with curvature at least $\delta\omega_Y$.
If $\mathcal J$ is its multiplier ideal, then
\begin{equation}\label{inj:fixed-ideal-inclusion}
 \mathcal O_Y(-D_*)\subseteq\mathcal J,
\end{equation}
since multiplying the metric coefficient by $|g_{D_*}|^2$
recovers the locally integrable density $F_\xi$.

\smallskip\noindent\textit{Interpolation and return to $X$.} Tensor by the bounded semipositive metric on $jkB$.  For every
$j\ge0$ the multiplier ideal remains the same $\mathcal J$,
which is not assumed trivial.  Nadel vanishing yields
\[
 H^v\bigl(Y,\mathcal O_Y(D_*+jkB)\otimes\mathcal J\bigr)=0
 \qquad(v>0,\ j\ge0).
\]
Put $\mathcal F=\mathcal O_Y(D_*)\otimes\mathcal J$.
Its Euler characteristic after twisting by $jkB$ is a polynomial in
$j$. The inclusion \eqref{inj:fixed-ideal-inclusion} gives
$\mathcal O_Y\hookrightarrow\mathcal F$, so at zero this polynomial
equals $h^0(Y,\mathcal F)>0$. A nonzero section therefore exists for
some $j>0$. Pull it back
and multiply by $\gamma^j$ to obtain an invariant section of
$jkP'+f^*D_*$.  Removing the exceptional twist and pushing down
gives the required section of $jkP$.
\end{proof}
